%!TEX root = ../hutbthesis_main.tex

\chapter{总结与展望}

\section{研究工作总结}

本课题围绕基于AirSim的无人机集群仿真软件性能优化展开，对多机仿真时出现的帧率急剧下降、时延大幅增加、CPU内核态占用过高、线程调度风暴等关键问题，进行了从理论分析、瓶颈定位、方案设计、源码改造直至实验验证的全流程系统研究。

主要完成工作如下：搭建标准化测试环境，重现并确认AirSim在20架以上无人机仿真时性能出现大幅度下滑的情况，运用专业性能分析工具找出根本原因，即线程模型不合理、sleep(0)低效等待、高频系统调用、调度风暴、内核态开销过大，提出并设计五种线程等待优化模型，对AirSim核心源码进行重构，实现线程池、负载均衡、CPU亲和性、单线程分发等关键优化，进行多维、大规模且长时间的实验，证实最优方案能让30架仿真时FPS从3提高到37，时延从292ms降至32ms，内核态占比从72\%降至22\%，上下文切换从286万次/秒降至8.3万次/秒，形成一套完整、可复用且可推广的高并发实时仿真系统性能调优范式，为机器人仿真、数字孪生、自动驾驶仿真等领域提供直接参考。。

\section{主要创新点}

（1）根因定位创新

首次将性能问题从``帧率低''的现象描述，深化到线程调度、内核态开销、超线程冲突的机理层面，完成精准瓶颈定位。

（2）优化模型创新

提出单线程分发器模型，用一个线程统一管理多机定时唤醒，从架构层面消灭调度风暴。

（3）工程方案创新

提供可配置、可切换、可对比的轻量化源码改造方案，不破坏原有功能，兼容全生态，易于落地使用。

（4）验证体系创新

构建``微观基准+集成仿真+长期稳定''三层验证体系，给出量化、可复现、高说服力的实验结论。

\section{研究不足与未来展望}

1.优化主要是针对单机平台展开的，并没有涉及到分布式集群仿真的扩展内容。未对GPU渲染、物理引擎并行计算进行进一步的优化。而且也没有在嵌入式平台（像NVIDIAJetson）上进行部署与测试工作。

2.未针对激光雷达、深度相机等高频传感器进行专门的数据收发优化。分布式多机协同仿真优化方面，把优化方案和分布式架构相结合，能够支持百架、千架级别的超大规模集群仿真。

3.GPU加速与异构计算融合上，借助CUDA加速物理更新、传感器数据处理，进而提高上限。嵌入式平台移植与部署是要实现在机载、边缘计算平台上的轻量级仿真，以此支持虚实融合测试。通用化性能调优框架是将优化思想封装成独立模块，能够快速移植到Gazebo、Webots、IsaacSim等平台。与集群智能算法深度结合则是提供专用性能接口，为编队控制、强化学习、协同感知等高端算法的研究提供服务。。
